\begin{theorem}[Complete metric limit uniqueness completion sibling route]
\label{thm:complete-metric-limit-uniqueness-completion-sibling-route}
Every $\CompleteMetricLimitUniquenessUp$ read that exports uniqueness must first consume the separated limit comparison supplied by \autoref{ch:concrete-instances-completemetric-namecert} and the completion-boundary row supplied by \autoref{ch:concrete-instances-metriccompletion-namecert}. The exported route therefore remains inside the displayed packet rows
$$
A,L_{0},L_{1},M,Z,S,H,C,P,N
$$
with $L_{0},L_{1}$ read as complete-metric limit rows over one Cauchy source, $M,Z,S$ read as the zero-distance and separated-classifier comparison, and $H,C,P,N$ preserving the same finite path. It exports no quotient completion point, selected representative, or host equality of metric points.
\end{theorem}
\begin{proof}
Project the carrier obligations of \autoref{thm:complete-metric-limit-uniqueness-namecert-obligations}. They expose the complete-metric limit rows $L_{0},L_{1}$ only over the shared source row $A$, so the separated comparison cannot begin from detached limit data.

Read the sibling completion boundary through \autoref{ch:concrete-instances-metriccompletion-namecert}. This supplies the completion-facing row consumed before $M,Z,S$ compare the two displayed limit reads, but it does not add quotient completion points or selected representatives.

Finally apply the separatedness handoff of \autoref{thm:complete-metric-limit-uniqueness-separatedness}. The support rows $H,C,P,N$ replay and name the same finite route, so the uniqueness export remains a completion-sibling dependency rather than host equality.
\end{proof}
